\begin{definitionbox}{Definition (Continuously Differentiable; Class \(C^1\))}
\begin{definition}[Continuously Differentiable; Class \(C^1\)]
\label{def:class-c1}
Let \(I\subseteq\mathbb{R}\) be an interval and let \(f:I\to\mathbb{R}\). The function \(f\) is \emph{continuously differentiable} on \(I\), written \(f\in C^1(I)\), if \(f\) is differentiable at every point of \(I\) and the derivative \(f':I\to\mathbb{R}\) is continuous on \(I\).
\end{definition}
\end{definitionbox}

\begin{remark*}[Standard quantified statement]
\[
  f\in C^1(I)
  \Longleftrightarrow
  \bigl(\forall x\in I:\operatorname{DifferentiableAt}(f,x)\bigr)
  \land
  \operatorname{ContinuousOn}(f',I).
\]
\end{remark*}

\begin{remark*}[Interpretation]
The condition \(f\in C^1(I)\) is strictly stronger than differentiability. A function may be differentiable everywhere while its derivative fails to be continuous; Darboux's Theorem still constrains that failure, because derivatives cannot jump. Thus the gap between differentiability and \(C^1\) is oscillatory rather than jump-like.
\end{remark*}

\begin{dependencies}
\begin{itemize}
  \item \hyperref[def:derivative-at-a-point]{Definition: Derivative at a Point}
  \item \hyperref[def:continuous-at-point]{Definition: Continuity at a Point}
\end{itemize}
\end{dependencies}

\begin{definitionbox}{Definition (Class \(C^k\))}
\begin{definition}[Class \(C^k\)]
\label{def:class-ck}
Let \(I\subseteq\mathbb{R}\) be an interval, let \(k\in\mathbb{N}\), and let \(f:I\to\mathbb{R}\). The function \(f\) is of \emph{class \(C^k\)} on \(I\), written \(f\in C^k(I)\), if the derivatives \(f^{(0)},f^{(1)},\ldots,f^{(k)}\) exist on \(I\) and \(f^{(k)}\) is continuous on \(I\). By convention, \(C^0(I)\) is the class of continuous functions on \(I\).
\end{definition}
\end{definitionbox}

\begin{remark*}[Standard quantified statement]
\[
  f\in C^k(I)
  \Longleftrightarrow
  \bigl(\forall j\in\{0,1,\ldots,k\}: f^{(j)}\text{ exists on }I\bigr)
  \land
  \operatorname{ContinuousOn}\bigl(f^{(k)},I\bigr).
\]
\end{remark*}

\begin{remark*}[Predicate reading]
\[
  \operatorname{ClassC}(f,I,k).
\]
For a fixed interval \(I\), the classes are nested:
\[
  f\in C^{k+1}(I)\Longrightarrow f\in C^k(I).
\]
\end{remark*}

\begin{remark*}[Interpretation]
The notation \(C^k(I)\) records \(k\) layers of differentiability together with continuity at the top layer. This is the regularity scale used to state how many times a theorem may differentiate a function without leaving the controlled setting.
\end{remark*}

\begin{dependencies}
\begin{itemize}
  \item \hyperref[def:higher-derivatives]{Definition: Higher Derivatives}
  \item \hyperref[def:class-c1]{Definition: Continuously Differentiable; Class \(C^1\)}
\end{itemize}
\end{dependencies}

\begin{definitionbox}{Definition (Smooth Function; Class \(C^\infty\))}
\begin{definition}[Smooth Function; Class \(C^\infty\)]
\label{def:class-cinfty}
Let \(I\subseteq\mathbb{R}\) be an interval and let \(f:I\to\mathbb{R}\). The function \(f\) is \emph{smooth}, or of \emph{class \(C^\infty\)}, written \(f\in C^\infty(I)\), if \(f\in C^k(I)\) for every \(k\in\mathbb{N}\). Equivalently,
\[
  C^\infty(I)=\bigcap_{k=0}^{\infty}C^k(I).
\]
\end{definition}
\end{definitionbox}

\begin{remark*}[Standard quantified statement]
\[
  f\in C^\infty(I)
  \Longleftrightarrow
  \forall k\in\mathbb{N}:\ f\in C^k(I).
\]
\end{remark*}

\begin{remark*}[Interpretation]
Smoothness is the regularity language imported by differential geometry. Smooth manifolds are assembled from \(C^\infty\) transition maps, and the tangent-map construction seeded by the differential uses this vocabulary.
\end{remark*}

\begin{dependencies}
\begin{itemize}
  \item \hyperref[def:class-ck]{Definition: Class \(C^k\)}
\end{itemize}
\end{dependencies}

\begin{definitionbox}{Definition (Real-Analytic Function; Class \(C^\omega\))}
\begin{definition}[Real-Analytic Function; Class \(C^\omega\)]
\label{def:class-comega}
Let \(I\subseteq\mathbb{R}\) be an interval and let \(f:I\to\mathbb{R}\). The function \(f\) is \emph{real-analytic}, or of \emph{class \(C^\omega\)}, written \(f\in C^\omega(I)\), if for every \(a\in I\) there exists \(r>0\) such that
\[
  f(x)=\sum_{n=0}^{\infty}\frac{f^{(n)}(a)}{n!}(x-a)^n
  \qquad\text{for every }x\in(a-r,a+r)\cap I.
\]
\end{definition}
\end{definitionbox}

\begin{remark*}[Standard quantified statement]
\[
  f\in C^\omega(I)
  \Longleftrightarrow
  \forall a\in I\;\exists r>0\;\forall x\in(a-r,a+r)\cap I:\;
  f(x)=\sum_{n=0}^{\infty}\frac{f^{(n)}(a)}{n!}(x-a)^n.
\]
\end{remark*}

\begin{remark*}[Interpretation]
Analytic regularity means that the Taylor series does more than approximate the function: near each point, the series equals the function. The infinite Taylor series and its convergence are developed in \hyperref[def:taylor-polynomial-at-a-point]{Taylor's construction and error analysis}, which supplies the series machinery behind this definition.
\end{remark*}

\begin{dependencies}
\begin{itemize}
  \item \hyperref[def:class-cinfty]{Definition: Smooth Function; Class \(C^\infty\)}
  \item \hyperref[def:taylor-polynomial-at-a-point]{Definition: Taylor Polynomial at a Point}
\end{itemize}
\end{dependencies}

\begin{theorembox}{Theorem (The Smoothness Tower)}
\begin{theorem}[The Smoothness Tower]
\label{thm:smoothness-tower}
For every nondegenerate interval \(I\), the inclusions
\[
  C^0(I)\supseteq C^1(I)\supseteq C^2(I)\supseteq \cdots
  \supseteq C^\infty(I)\supseteq C^\omega(I)
\]
hold, and each displayed inclusion is strict.
\hyperref[prf:smoothness-tower]{Go to proof.}
\end{theorem}
\end{theorembox}

\begin{remark*}[Standard quantified statement]
\[
  \forall k\in\mathbb{N}:\ C^{k+1}(I)\subsetneq C^k(I),
  \qquad
  C^\omega(I)\subsetneq C^\infty(I).
\]
\end{remark*}

\begin{remark*}[Interpretation]
The inclusions are definitional; strictness is witnessed by examples. On any nondegenerate interval, translated and rescaled copies of \(x\mapsto x^k|x|\) separate \(C^k\) from \(C^{k+1}\). Translated copies of the flat function
\[
  f(x)=
  \begin{cases}
    e^{-1/x^2}, & x\neq 0,\\
    0, & x=0,
  \end{cases}
\]
separate \(C^\infty\) from \(C^\omega\): all derivatives at the flat point vanish, so the Taylor series is identically zero, while the function is positive on one side of that point.
\end{remark*}

\begin{dependencies}
\begin{itemize}
  \item \hyperref[def:class-ck]{Definition: Class \(C^k\)}
  \item \hyperref[def:class-cinfty]{Definition: Smooth Function; Class \(C^\infty\)}
  \item \hyperref[def:class-comega]{Definition: Real-Analytic Function; Class \(C^\omega\)}
  \item \hyperref[thm:darboux]{Theorem: Darboux's Theorem}
\end{itemize}
\end{dependencies}
